\chapter{Contexte, problématique et démarche}
\chapminitoc

\section{Contexte : l'existant à mon arrivée}
% TODO : état des deux véhicules à ton arrivée (ce qui marchait, ce qui
% manquait), comme le « Contexte » de tes chapitres 4A.

\section{Problématique}
% La question unique qui chapeaute tout le rapport — encadrée, bien visible.
\begin{center}
\fbox{\parbox{0.9\textwidth}{\itshape
Comment transformer des prototypes de véhicules à échelle réduite en une
plateforme téléopérable, autonome et communicante --- fiable du niveau
mécanique jusqu'au niveau applicatif --- tout en instrumentant ses performances
réseau (WiFi/5G) pour en démontrer la viabilité ?}}
\end{center}
% TODO : reformule-la avec tes mots, puis décline-la en sous-questions.

\section{Objectifs et cahier des charges}
% TODO : tableau objectifs / livrables / critères de réussite
% (téléop 5G opérationnelle, V2V CAM+DENM démontré, base Protova1 fiabilisée,
% dashboard fonctionnel, etc.)

\section{Démarche et planning}
% TODO : Gantt prévisionnel vs réalisé (figure), méthode de travail
% (itérations, démos régulières).

\section{Moyens mis à disposition}
% TODO : matériel (Jetson, RPLIDAR A1, caméra Orbbec, dongle 5G, volant G29,
% imprimante 3D...) et logiciels (ROS 2 Humble, Inventor/CAO, Vanetza...).

\section{Organisation et communication}
% TODO : rituels d'équipe, échanges avec les tuteurs (CDC envoyé, bilans
% intermédiaires) — la grille INSA note explicitement ce point.
